% =====================================================================
% Appendix: Test log
% Sources: LM-cooling/Experimental Data/Experimental Measurement.xlsx and
% Thermistor_data/2 DegC values (CSV numbers as in Thermistor_data/TOC.txt)
% =====================================================================
\chapter{Test Log}
\label{app:test-log}

\Cref{tab:app-test-log} lists every thermal test run on the two loops in
chronological order, with the data file that holds its temperature record and
where its results are used in this report. The two loops were usually run side
by side, starting together, so the Galinstan and water tests of one day share an
ambient temperature. The thermistor logger recorded every temperature
continuously, about once every \SI{0.5}{\second}; the steady-state values in
\cref{tab:measured-performance,tab:tn_results} are the K-type thermocouple
readings recorded near the end of each heated period, together with the heater
power, flow and pump power, in \texttt{Experimental Measurement.xlsx}.

\begin{table}[htbp]
  \centering
  \caption{Log of the thermal tests.}
  \label{tab:app-test-log}
  \footnotesize
  \setlength{\tabcolsep}{3pt}
  \begin{tabularx}{\linewidth}{@{}l c >{\raggedright\arraybackslash}p{3.6cm}
      c c c c c L@{}}
    \toprule
    Date & CSV & Test & $\Vdot$ & $\Qin$ & $T_{\mathrm{A}}$ & Logged & Steady
      & Used in \\
         &     &      & (\si{\milli\litre\per\second}) & (\si{\watt})
      & (\si{\celsius}) & (\si{\minute}) & (\si{\minute}) & \\
    \midrule
    2026-07-28 & 6      & Water, power sweep        & 59--62 & 10--90 & 22.3 & 283 & --  & Preliminary \\
    2026-07-31 & 7      & Water, steady state       & 29.2   & 72     & 22.3 & 122 & --  & Preliminary \\
    2026-07-31 & 9      & Water, step to \SI{110}{\watt} & 30.8 & 70--110 & 22.7 & 96 & -- & Water pump COP\tnote{a} \\
    2026-08-02 & 10     & Galinstan, no flow        & 0      & 12     & 21.2 & 207 & 171 & $\Rmisc$ \\
    2026-08-02 & 11     & Water, no flow            & 0      & 12     & 21.2 & 207 & 169 & $\Rmisc$ \\
    2026-08-02 & 12--15 & Both, propagation and cool-down & --  & --  & 21.2 & --  & --  & Not used \\
    2026-08-05 & 16     & Galinstan, full flow (\SI{18}{\ampere}) & 1.45 & 72 & 20.4 & 161 & 50 & \Cref{tab:measured-performance,tab:tn_results}; \cref{fig:galinstan_transient}(a) \\
    2026-08-05 & 17     & Water, full flow (\SI{100}{\percent} duty) & 58.3 & 72 & 20.4 & 160 & 61 & \Cref{tab:measured-performance,tab:tn_results}; \cref{fig:water_transient}(a) \\
    2026-08-06 & 18     & Galinstan, low flow (\SI{9}{\ampere}) & 0.75 & 72 & 20.3 & 64 & -- & Stopped at the \SI{70}{\celsius} pump limit \\
    2026-08-06 & 19     & Water, low flow (\SI{20}{\percent} duty) & 29.2 & 72 & 20.3 & 184 & 77 & \Cref{tab:measured-performance,tab:tn_results}; \cref{fig:water_transient}(b) \\
    2026-08-10 & 20     & Galinstan, low flow (\SI{9}{\ampere}) & 0.77 & 40 & 21.9 & 105 & 52 & \Cref{tab:measured-performance,tab:tn_results}; \cref{fig:galinstan_transient}(b) \\
    2026-08-10 & 21     & Water, low flow (\SI{20}{\percent} duty) & 29.2 & 40 & 21.9 & 105 & -- & Not used \\
    2026-09-07 & 24     & Water, restricted (\SI{0.73}{\litre\per\minute}) & 12.2 & 72 & 22.8 & 103 & 48 & \Cref{tab:measured-performance,tab:tn_results}; \cref{fig:water_transient}(c) \\
    2026-09-08 & 25     & Water, restricted (\SI{0.5}{\litre\per\minute}) & 8.3 & 72 & 21.7 & 170 & 94 & \Cref{tab:measured-performance,tab:tn_results}; \cref{fig:water_transient}(d) \\
    2026-09-15 & 29     & Resistors, no cooling system & --  & 12  & 22.6 & 239 & --  & Not used \\
    2026-09-16 & 31     & Resistors, no cooling system & --  & 12  & 22.5 & 254 & --  & Not used \\
    2026-09-24 & 34, 35 & Both, small resistor      & 8.3, 1.37 & 25 & 21.3 & 108 & -- & Not used \\
    2026-10-05 & 36     & Galinstan, inductor (\SI{18}{\ampere}) & 1.4 & 42.8\tnote{b} & 21.8 & 231 & 143 & Experiment~2; \cref{fig:inductor-cam} \\
    \bottomrule
  \end{tabularx}
  \floatnote{CSV: file number in \texttt{Thermistor\_data/2 DegC values}.
    Logged: length of the record. Steady: time after which the resistor (or
    inductor) temperature, as a \SI{2}{\minute} moving average, stayed within
    \SI{0.5}{\celsius} of its final value until the heat was switched off.
    \tnote{a}\Cref{tab:pumps}. \tnote{b}Estimated loss
    (\cref{sec:model-validation}).}
\end{table}
\issue{Confirm that CSV~9 (step from \SI{70}{\watt} to \SI{110}{\watt}) is the
additional \SI{110}{\watt} test behind the water COP of 47.6 in \cref{tab:pumps}.}

Each steady-state test reached steady state within \SIrange{48}{94}{\minute} of
the heat being switched on, and was held there for a further
\SIrange{37}{94}{\minute} before the heat was switched off. The no-flow tests
took about \SI{170}{\minute}, because without flow the whole apparatus heats
slowly by conduction. The inductor test took about \SI{140}{\minute}, as the
upper core, which is cooled mainly by natural convection, has a long thermal time
constant.

The project files that hold these results are:
\begin{itemize}
  \item \texttt{Experimental Data/Experimental Measurement.xlsx}: thermocouple
    readings, heater power, flow and pump power at steady state for each test,
    and the Galinstan pump and inductor measurements;
  \item \texttt{Experimental Data/Thermistor\_data/2 DegC values}: the
    temperature records of \cref{tab:app-test-log}, with a table of contents in
    \texttt{Thermistor\_data/TOC.txt};
  \item \texttt{Experimental Data/Inductor}: the thermal camera images and scope
    captures of Experiment~2; and
  \item \texttt{Final Report/Model Appendix/Thermal\_Network\_Model.xlsx}: the
    thermal network of \cref{sec:thermal-network}, the $\varepsilon$--NTU
    calculations of \cref{app:convection} and the predictions of
    \cref{tab:tn_results}.
\end{itemize}
